\section{Introduction}
\label{sec:intro}
\textcolor{red}{currently organize into subsections, then later connect those subsections into a single thread}

Large language models (LLMs)~\cite{wei2022emergentabilitieslargelanguage, openai2024gpt4technicalreport,deepseekai2025deepseekv3technicalreport,wei2022emergent} have been increasingly deployed across various domains, achieving remarkable performance in applications including conversational assistants~\cite{dam2024completesurveyllmbasedai,10.1145/3580305.3599572}, code generation~\cite{hui2024qwen25codertechnicalreport}, and retrieval-augmented generation (RAG)~\cite{NEURIPS2020_6b493230,yao2025cacheblendfastlargelanguage}. They have started to profoundly impact our work and daily routines. Many cloud companies~\cite{openai2024chatgpt,anthropic2024claude,google2024gemini} are racing to provide these applications as hosted services. However, running these applications is very expensive, requiring a large number of hardware accelerators such as GPUs. As deployment grows, the latency and resource efficiency of both online and offline inference have become increasingly critical. In response, a plethora of systems have been deployed to optimize different aspects of LLM serving, such as scheduling strategies \cite{pang2025optimizingllminferencethroughput,zhong2024distservedisaggregatingprefilldecoding,zheng2025bucketservebucketbaseddynamicbatching,agrawal2024tamingthroughputlatencytradeoffllm}, parallelization techniques \cite{su2025seesawhighthroughputllminference,wu2024loongserveefficientlyservinglongcontext}, memory management \cite{kwon2023efficientmemorymanagementlarge,prabhu2025vattentiondynamicmemorymanagement,zhang2025jengaeffectivememorymanagement}, quantization \cite{chee2024quip2bitquantizationlarge,lin2024awqactivationawareweightquantization}, and kernel-level performance \cite{yi2025patacceleratingllmdecoding,MLSYS2025_96894468}.

To understand the landscape of optimizations covered in this paper, we briefly review the key technical foundations. LLM inference operates in two distinct phases: a compute-bound \textbf{prefill} phase that processes all prompt tokens in parallel to produce the first output token, and a memory-bound \textbf{decode} phase that generates one token per iteration. The decode phase relies on a \textbf{KV cache} — stored keys and values from prior tokens — which grows with context length and increases the memory footprint.

At the hardware level, GPUs execute thousands of parallel threads across Streaming Multiprocessors (SMs), each with local shared memory and tensor cores, all backed by a shared L2 cache and global memory; the interplay between compute throughput and memory bandwidth at this level shapes the design of nearly every optimization discussed in this paper. 

These characteristics of LLM inference motivate a large body of work on \textbf{batching and scheduling}. Because prefill and decode have mismatched compute profiles, systems have explored strategies that either  interleave~\cite{agrawal2024tamingthroughputlatencytradeoffllm}, disaggregate~\cite{zhong2024distservedisaggregatingprefilldecoding}, or spatially co-schedule~\cite{Kamath_2025,lin2025boostingllmservingspatialtemporal} the two phases. 

Many practical workloads, such as few shot prompting~\cite{brown2020languagemodelsfewshotlearners,NEURIPS2021_5c049256}, system prompts~\cite{thebigpromptlibrary}, and RAG~\cite{lewis2021retrievalaugmentedgenerationknowledgeintensivenlp}, also exhibit shared prompt prefixes across requests. Reusing the KV cache of these shared prefixes avoids redundant computation and reduces memory traffic. Systems such as SGLang~\cite{zheng2024sglangefficientexecutionstructured} and vLLM~\cite{kwon2023efficient} manage this through radix trees and Merkle-tree-inspired hashing schemes, respectively. 

A complementary line of kernel-level work~\cite{yi2025patacceleratingllmdecoding,yao2025deftdecodingflashtreeattention,ye2024chunkattentionefficientselfattentionprefixaware,MLSYS2025_96894468} pushes this idea further by restructuring how queries are grouped and KV data is loaded within the attention kernel itself, enabling the reuse of shared KV data and reducing redundant global memory traffic at the GPU execution level.

Despite these advances, a precise empirical understanding of \textit{why} and \textit{how} prefix structure affects decode throughput has been lacking. In this paper, we present a series of controlled experiments that systematically quantify the relationship between prefix structure and inference performance.

Our central finding is that decode-phase throughput is primarily governed by KV cache memory access patterns rather than raw arithmetic compute. (\textcolor{red}{should we reword the previous statement?}) Intuitively, during the decode phase, every new token generated by a request requires attending to \textit{all} previously seen tokens, which means that the GPU must repeatedly read through the entire KV cache at every generation step. When multiple requests in a batch share the same prompt prefix, they traverse the same regions of GPU memory at nearly the same time. This synchronized access improves both \textit{spatial locality} (the hardware prefetcher can reliably anticipate and load upcoming memory regions) and \textit{temporal locality} (a memory block fetched for one request is likely to be reused by a neighbouring request before being evicted). Together, these effects increase L2 cache hit rates and sustain high memory bandwidth utilization. 

We validate this hypothesis through several controlled experiments. First, we show that even a single request, out of a large batch, with a different prefix, can reduce decode throughput by nearly $2\times$. This stark sensitivity arises because the GPU memory access stream becomes interleaved across two disjoint KV cache regions, disrupting prefetcher confidence and reducing temporal cache reuse. Second, we demonstrate that throughput improves smoothly and monotonically as the fraction of shared prefix tokens increases: moving from no sharing to full sharing yields over 60\% higher throughput. Third, we find that the \textit{position} of sharing within a sequence also matters; since all requests traverse the KV cache in lockstep during decode, sharing at the beginning of a sequence provides the greatest locality benefit. Once requests diverge, synchronized memory access is disrupted, and subsequent sharing in later parts of the sequence provides insignificant returns. Fourth, beyond locality benefits, we show that when the number of simultaneously active prefixes becomes large enough to exhaust GPU memory capacity, the system is forced to evict and recompute KV cache blocks, causing an additional sharp decline in throughput and effective batch size. This highlights that prefix-aware scheduling must account for both memory access locality and GPU memory capacity constraints. \textcolor{red}{reword the last sentence, maybe does not fit our narrative}

We further characterize the interplay between batch size and prefix structure. Increasing batch size initially improves throughput by better utilizing GPU tensor cores; however, once compute resources are saturated, adding more requests only increases KV cache traffic without proportional throughput gains, making memory bandwidth the bottleneck. Under end-to-end execution, including both prefill and decode phases, we show that prefix-homogeneous batches can outperform larger heterogeneous batches once decode lengths are sufficiently long. This is because each decode step requires another full sweep of the KV cache, allowing the locality benefits of aligned memory access to compound over time. Interestingly, even a homogeneous batch of 50 requests can surpass a heterogeneous batch of 800 requests when decode lengths are long enough, demonstrating that memory locality can outweigh the loss in tensor core utilization.

Finally, we examine the cost of detecting prefix sharing at runtime. Existing systems such as SGLang~\cite{zheng2024sglangefficientexecutionstructured} use radix-tree-based schedulers that perform token-by-token comparisons, dynamic node splitting, and repeated tree traversals to identify shared prefixes. We show that for a workload of $2000$ requests with 10k-token prefixes, these scheduling policies incur CPU overhead of 78-89 seconds, which is on the same order as the GPU execution time itself. In contrast, when prefix group identities are provided explicitly by the user, the scheduling overhead drops to just 0.03 seconds (\textcolor{red}{no need to specify the ideal policy here, rename by oracle or something}). This reveals a fundamental inefficiency: the very mechanism used to exploit prefix sharing can itself become a bottleneck, motivating a shift towards a lightweight prefix identification scheme.

\textcolor{red}{should the metric come out here?}

Motivated by these insights, we propose a scheduling system built around two complementary components. The first is the \textbf{Chunked Hash Tree (CHT)}, a lightweight data structure for fast prefix detection. Rather than token-by-token radix tree traversals, CHT represents each request as a compact vector of cumulative prefix hashes computed at fixed token intervals, enabling $O(\log W)$ identification of the best prefix-sharing candidate among $W$ waiting requests. A min-heap ordered by per-request \emph{missing counts} — the number of prefix chunks not yet covered by the active batch — allows the scheduler to efficiently surface the request that maximizes prefix locality, while lazy heap updates keep CPU overhead minimal, even for thousands of long-prompt requests.

The second component addresses the question of \emph{when to stop} adding requests to a batch. Growing a batch improves GPU utilization, but each new request with a diverging prefix shrinks the overall shared prefix length and reduces locality. We frame this stop-or-continue decision as a reinforcement learning problem. At each step of batch construction, the scheduler observes the current batch size, the shared prefix length before and after admitting the next candidate, and the number of waiting requests that share the same prefix level (\textcolor{red}{info about all the state variables not needed probably}). Tabular policies, trained online from observed decode throughput, learn to balance these competing objectives adaptively across different GPU configurations and workload conditions, which a fixed heuristic policy cannot do reliably. Together, the CHT and the learned batching policy form an end-to-end scheduling system that simultaneously reduces CPU scheduling overhead and improves decode-phase throughput by constructing large, prefix-homogeneous batches. 


\subsection{\textcolor{blue}{Key conclusions/insights from these challenges}}

\subsection{\textcolor{blue}{Contributions of this paper}}
In summary, the paper makes the following contributions:
\begin{itemize}
    \item We conduct an in-depth analysis of the bottlenecks that lead to reduced DRAM bandwidth utilization for prefix-shared workloads.
    \item We also identify several CPU-level overheads while detecting prefix sharing in a batch of requests.
    \item We design and implement an efficient data structure that effectively organizes requests and reduces the scheduler's CPU overhead for such workloads.
    \item We come up with a metric that, in turn, drives our batching decisions. We also employ reinforcement learning policies to guide batch formation.
    \item Experimental evaluations on real-world workloads show that we outperform state-of-the-art inference engines in both latency and throughput while achieving significantly better DRAM bandwidth utilization.
\end{itemize}

\subsection{\textcolor{blue}{Overall paper organization}}
We present related work in this field in \S\ref{sec:background}. We provide an extensive set of experiments that demonstrate our insights and key ideas in \S\ref{sec:motivation}. We next present the design and implementation of our ideas in \S\ref{sec:design}. We evaluate our implementation with real workloads against state-of-the-art baselines in \S\ref{sec:eval}, along with micro-benchmarks and an ablation study, followed by the conclusion in \S\ref{sec:conclusion}.